\documentclass[runningheads]{llncs}

\usepackage[T1]{fontenc}
\usepackage{graphicx}
\usepackage{amsmath}
\usepackage[hidelinks]{hyperref}

\begin{document}

\title{Beyond Student Labels: Evaluating Prompt Sensitivity in AI-Generated Practice Problems for Multilingual Learners}
\titlerunning{Prompt Sensitivity in AI-Generated Practice Problems}

\author{Satyanarayana Rudraraju}
\authorrunning{S. Rudraraju}

\institute{Whitney M. Young Magnet High School, Chicago, IL, USA \\ \email{rudrarajumc@gmail.com}}

\maketitle

\begin{abstract}
When I started tutoring multilingual students in Hyderabad, India, most of whom speak Telugu at home while studying math in English, I noticed something that bothered me. If I told an AI model that a student was learning English, the problems it wrote looked easier, not just in the wording but in the math itself. It could have been a fluke or a real pattern affecting my students. This paper is my attempt to figure out which. I ran a pre-registered experiment testing practice-problem writing across three ability levels, four language descriptions, three prompt rewordings, five middle-school topics, and three repeats, on three free, open-weight models (Llama 3.3 70B, GPT-OSS-120B, Mistral Small), collecting all 1,620 generations. The main finding is a null: after Holm-Bonferroni correction, no language label significantly changed math difficulty (pooled Cohen's $d \leq 0.25$), and the null held after removing near-duplicate outputs, clustering standard errors, and restricting to English-only output. The models did simplify language, though: ELL and multilingual-learner labels lowered reading grade level by about 0.9 grades (corrected $p<.05$), appropriate adaptation to English proficiency, not a proxy for math ability. Three further findings emerged that a solve-only benchmark would miss: models unpromptedly switched languages, writing 53.7\% of ELL-Spanish problems entirely in Spanish and 19.5\% of ELL-Mandarin problems in Chinese; character names strongly tracked the label ($\chi^2(9)=365.6$, $p<10^{-72}$); and an AI judge scored all three problems I flagged for cultural bias a perfect 5/5, while a human reader caught them. Code, data, and the pre-registration are public.

\keywords{Large language models \and Algorithmic fairness \and Multilingual learners \and Educational technology \and Prompt sensitivity \and LLM-as-judge \and Pre-registration.}
\end{abstract}

\section{Introduction}

This study grew out of a free tutoring program I run in Hyderabad, India, for multilingual students who study math in English while speaking Telugu at home. When I told a large language model that a student was learning English, the practice problems it generated looked easier, not just in wording, but in the underlying mathematics. This paper is my attempt to find out whether that is a real, systematic pattern.

The concern has roots in education research. Low teacher expectations, even unspoken, can reduce student performance \cite{rosenthal1968,steele1995}, and English-language-learner (ELL) classification specifically has been shown to lower teachers' academic expectations independent of ability \cite{umansky2021}; tracking research documents a similar dynamic with curricular material \cite{oakes1985}. An ELL label describes English proficiency, not mathematical ability. An AI system that treats the two as correlated would automate a documented human bias at the scale of every generated problem.

I study problem writing, not problem solving. Prior multilingual math benchmarks test whether models can solve translated problems \cite{cobbe2021,shi2022}; far less is known about what models produce when asked to write practice material for a described student, even as AI-generated problem writing becomes a standard tutoring-tool feature \cite{kasneci2023}. Writing is where a label has room to leak into content: the model chooses the numbers, the operations, the names, the setting, and (it turns out) sometimes the language itself, none of which a fixed-problem benchmark can observe.

I deliberately test free, open-weight, no-credit-card models, not the top paid systems, but the models a low-budget tutoring program or under-resourced school can actually run. If a fairness problem exists at this tier, it is the tier the least-supported students will meet.

This paper contributes: (1) a pre-registered design separating math difficulty from language difficulty under an ELL label; (2) a test of robustness across three independent prompt rewordings; (3) a measurement of unprompted first-language switching, a behavior only observable in a writing task; and (4) evidence of a cultural blind spot in LLM-as-judge scoring relative to human review.

\section{Related Work}

\textbf{Persona and label bias.} Assigning a model a described identity changes its outputs and can degrade performance for some groups \cite{gupta2023}; broader audits find demographic details push generations toward stereotypes \cite{sheng2019,deshpande2023}, with effects varying enough by provider that each model needs independent audit \cite{liang2022,gallegos2024}. Following Blodgett et al.'s call for an explicit normative grounding in bias claims \cite{blodgett2020}, mine is stated directly: a language label should change language, not mathematics. Most prior audits examine opinions, refusals, or answer accuracy under a persona; I ask a narrower, school-specific question: under a language-learner identity, does the model change the math or only the words?

\textbf{AI-generated educational content.} MATHWELL \cite{christ2024} evaluates teacher-annotated, model-written word problems and finds models often miss the requested grade level, implying difficulty is something models actively (and sometimes unreliably) control, which is the premise my manipulation check (H1) tests.

\textbf{Prompt sensitivity.} Open models are highly sensitive to meaning-preserving prompt reformatting \cite{sclar2023}, motivating my robustness check across three independent wordings (H4).

\textbf{LLM-as-judge.} Cross-model judges show high aggregate agreement with humans but documented biases toward position, length, and same-family output \cite{zheng2023}. I use cross-family judging and a human-blinded validation sample, and report a specific judge-human disagreement on cultural content (Sect.~4.5) as a finding in its own right.

No prior study I found isolates the math-versus-language trade-off under an ELL label with pre-registered, cross-prompt robustness, measures unprompted language-switching as a generation behavior, or tests whether an LLM judge detects human-flagged cultural issues.

\section{Method}

\subsection{Research Questions and Design}

Five hypotheses were pre-registered (\texttt{DESIGN.md}, 2026-07-03) before data collection: H1 (manipulation check) math difficulty rises with stated ability level; H2 (headline) ELL labels lower math difficulty versus a no-mention control; H3 (adaptation) ELL labels lower language difficulty; H4 (robustness) effects are stable across reworded prompts; H5 (cultural content) language-specific labels change the names and settings used. The logic separates fair adaptation (H3 true, H2 false) from the documented harm of treating a language label as a math-ability signal (H2 true).

The design is a fully crossed factorial (Table~\ref{tab:design}): 540 problems per model family, times three families, giving \textbf{N = 1{,}620} generations, all collected. Grade level (7th) and temperature (0.7) were fixed. Pooled power is adequate ($\approx.99$ for $d=0.3$, $\approx.80$ for $d=0.2$); a single within-model cell ($n\approx130$) reaches only $\approx.67$ power for $d=0.3$, reaching 80\% power only around $d\approx0.35$, a correction to the pre-registration's overstated power claim, disclosed here rather than silently fixed.

\begin{table}[t]
\caption{Factorial design (\texttt{DESIGN.md} \S3; \texttt{src/build\_prompt\_matrix.py}).}
\label{tab:design}
\centering\small
\begin{tabular}{p{2.6cm}p{7.7cm}}
\hline
Factor & Levels \\
\hline
Stated ability (3) & struggling, on-level, advanced \\
Language description (4) & control, multilingual, ELL-Spanish, ELL-Mandarin \\
Prompt wording (3) & Templates A, B, C (same meaning, different phrasing) \\
Topic (5) & proportions, percents, two-step equations, rational ops, circles \\
Repeat (3) & independent samples, temperature 0.7 \\
\hline
\end{tabular}
\end{table}

Language description is a short clause appended to an otherwise identical sentence; the control condition appends an empty string. For example (topic = percents, level = struggling), the control reads ``...for a 7th-grade student who is struggling with this topic;'' the ELL-Spanish condition appends ``...and is an English language learner whose first language is Spanish.'' Three independently worded templates (A/B/C) express the same request in different phrasing; full text is in the public repository.

Three pinned, free, no-credit-card model families were used: Llama 3.3 70B (Meta, via Groq), GPT-OSS-120B (OpenAI's open-weight model, via Cerebras), and Mistral Small (Mistral AI); exact pinned API identifiers are recorded in \texttt{config.json} in the repository. A fourth family, Gemini 2.5 Flash (with 24 of its responses served by Gemini 2.5 Flash-Lite), is reported descriptively only (Sect.~5) after its free tier returned too few usable items. Data were collected 2026-07-03 to 2026-07-07. These models (24B-120B parameters) represent an access tier, not a state-of-the-art tier, by design.

\subsection{Measures and Judging}

Math difficulty: solution-step count, operation count, and largest number used, read from the model's own solution. Language difficulty: Flesch-Kincaid grade \cite{kincaid1975}, sentence length, and rare-word ratio \cite{speer2022}, computed on English-language output only. A cross-family LLM judge scored five 1-5 dimensions (correctness, difficulty fit, clarity, pedagogical soundness, cultural neutrality) and extracted names and settings for the H5 analysis.

Each readable problem was scored by a judge from a different model family than its writer, blind to condition labels, in a fixed cross-family rotation. I independently and blindly rated a stratified 56-item sample (batch 1, covering GPT-OSS and Mistral) against a pre-registered pass mark of Spearman $\rho \geq 0.5$ on correctness.

\subsection{AI Tool Use Disclosure}

This research used Claude (Anthropic) to assist with manuscript drafting, statistical code review, and visualization generation. I performed all substantive scholarly work myself (design, hypotheses, data collection, human rating, analysis, and interpretation) and independently verified every reported statistic against the raw data. No figure is AI-generated: all figures are deterministic \texttt{matplotlib} output regenerable from the released code, and all table values are computed by the released analysis script. The 1,620 practice problems and the judge scores are themselves produced by the language models under audit, which is the object of the study, not a shortcut in producing it.

\subsection{Analysis and Protocol Amendments}

The primary test regresses each outcome on ability level $\times$ language, prompt wording, topic, and model family, with HC3 robust standard errors \cite{mackinnon1985}, reporting the effect versus control, its 95\% CI, and Cohen's $d$ \cite{cohen1988}, with Holm-Bonferroni correction \cite{holm1979} applied within each measure family (math / language / judge) at $\alpha=.05$. H5 uses a chi-square test on name origin by condition.

Four protocol changes were made before hypothesis testing and are disclosed with timestamps in \texttt{DESIGN.md}: a provider retirement forced a Qwen$\to$GPT-OSS substitution before any data existed in that slot; judge rotation was reassigned for free-tier quota; output-language detection was added via \texttt{langdetect} \cite{nakatani2010} after observing first-language switching in matched pairs; and Gemini was dropped from confirmatory analysis after its free tier returned only 29 of 540 planned usable items.

\section{Results}

Of 1,620 generations, 1,567 (96.7\%) produced readable JSON; parse failure was independent of language condition ($\chi^2=3.34$, $p=.34$). Table~\ref{tab:descriptives} gives condition means.

\begin{table}[t]
\caption{Mean outcomes by condition (readable items). Math columns: $n=1{,}567$. FK grade: English outputs only, $n=1{,}121$. Sentence length: $n=1{,}496$.}
\label{tab:descriptives}
\centering\small
\begin{tabular}{lccccc}
\hline
Condition & Steps & Ops & MaxNum & FK grade & Sent.\ len.\\
\hline
Control & 5.07 & 17.93 & 67.55 & 4.64 & 8.13 \\
Multilingual & 5.00 & 16.38 & 67.40 & 4.25 & 8.03 \\
ELL-Spanish & 4.57 & 15.44 & 47.06 & 4.22 & 7.85 \\
ELL-Mandarin & 4.46 & 13.88 & 56.82 & 3.91 & 7.20 \\
\hline
\end{tabular}
\end{table}

\subsection{H1: Manipulation Check}

An advanced-level label raised solution steps by 1.84 ($p<.001$), operations by 10.28 ($p<.001$), and the largest number by 120.8 ($p<.001$) versus on-level, corroborated by judge-scored difficulty fit and correctness drops (both $p<.001$): advanced problems are genuinely harder, not just longer. The struggling label, however, did not reliably lower any math measure (steps $p=.09$; operations $p=.60$; largest number $p=.41$): models raised difficulty on request but resisted lowering it. This asymmetry means the design is better powered to detect a difficulty increase than a decrease, exactly the direction H2's hypothesized harm would take, and bounds how strong the eventual null can be read.

\subsection{H2: No Significant Change in Math Difficulty}

No language label significantly changed any math measure after Holm-Bonferroni correction (all corrected $p \geq .24$). Table~\ref{tab:permodel} gives per-model Cohen's $d$ for the solution-steps effect; pooled effects are small ($d=-0.02$ to $-0.25$) with confidence intervals crossing zero. Per-model estimates vary, with Llama's ELL-Spanish effect ($d=-0.42$) and Mistral's ELL-Mandarin effect ($d=-0.38$) the largest, versus GPT-OSS's largest effect of $-0.13$ (Fig.~\ref{fig:forest}), but these comparisons are exploratory and uncorrected. The pre-registered falsification clause specified a null as $|d|<0.2$ with crossing CIs; the pooled ELL-Mandarin estimate ($-0.25$) sits slightly outside that band, so I report a null with a small residual trend rather than a claim of exactly zero effect.

\begin{table}[t]
\caption{Cohen's $d$ for the language effect on solution steps, by model (exploratory, uncorrected).}
\label{tab:permodel}
\centering\small
\begin{tabular}{lccc}
\hline
Model & Multilingual & ELL-Spanish & ELL-Mandarin \\
\hline
Llama 3.3 70B & $-0.06$ & $-0.42$ & $-0.29$ \\
Mistral Small & $-0.02$ & $-0.23$ & $-0.38$ \\
GPT-OSS-120B & $-0.03$ & $-0.10$ & $-0.13$ \\
\textbf{Pooled} & $-0.02$ & $-0.19$ & $-0.25$ \\
\hline
\end{tabular}
\end{table}

\begin{figure}[t]
\centering
\includegraphics[width=0.85\linewidth]{figures/effect_sizes_forest.png}
\caption{Effect of the language label on math difficulty (solution steps), by model and pooled. The shaded band marks a small effect ($|d|<0.2$). No effect is significant after correction.}
\label{fig:forest}
\end{figure}

\subsection{H3: Language Adaptation}

The same models that left math unchanged significantly simplified language. On Flesch-Kincaid grade (English outputs, $n=1{,}121$), the multilingual label lowered reading grade by 0.89 (corrected $p=.0071$) and ELL-Mandarin by 0.92 (corrected $p=.0120$); ELL-Spanish trended the same direction ($-0.55$) but fell short after correction ($p=.086$ uncorrected; Fig.~\ref{fig:fk}). One rare-word effect survived correction (ELL-Mandarin $+0.020$, $p=.0415$), but this reflects a naming artifact (frequency counters score romanized Chinese names as rare English words) rather than harder vocabulary, confirmed by the name analysis in Sect.~4.4. Set against H2's null, the one-line finding is: the models adapted the language, not the math.

\begin{figure}[t]
\centering
\includegraphics[width=0.85\linewidth]{figures/fk_grade_by_language.png}
\caption{Flesch-Kincaid reading grade by language condition, one line per model, with standard-error bars.}
\label{fig:fk}
\end{figure}

\subsection{Robustness Checks}

Three robustness checks confirm the null is not an artifact. Removing near-duplicate outputs (107 of 1,567 items, $\geq 90\%$ similarity within cell) left the null intact (smallest corrected $p=.70$); Llama's ELL-Spanish effect grew, not shrank, after de-duplication ($d=-0.42\to-0.51$), ruling out repetition as the driver. Clustering standard errors on the 539 design cells weakened the strongest uncorrected effect ($p=.07\to.11$). Restricting to English-only output, which rules out a parsing artifact, preserved the null (smallest uncorrected $p=.12$). Separately, the ELL-versus-control gap kept a consistent negative sign across all three prompt wordings for every math measure (H4), though its magnitude varied.

\subsection{H5: Cultural Content}

Character names tracked the label strongly ($\chi^2(9)=365.6$, $p<10^{-72}$, Cram\'er's $V=0.28$; Table~\ref{tab:names}, Fig.~\ref{fig:names}). Chinese-style names appeared in 15.9\% of ELL-Mandarin problems and 0.0\% elsewhere ($\chi^2=191.3$, $p<10^{-42}$); Hispanic-style names appeared in 21.8\% of ELL-Spanish problems versus 6.3\% elsewhere ($\chi^2=75.3$, $p<10^{-17}$). The generic ``multilingual learner'' label also pulled toward Hispanic names (72 items, mostly ``Maria''), despite naming no specific language. Whether this matching constitutes harm is debatable: it can read as respectful localization or as a narrow, stereotyped default. What is not debatable is the second half of H5: in the 56-item human check, three problems were flagged for cultural stereotyping (two Spanish-language garden/circle problems, one Mandarin-condition naming choice), and the cross-family judge scored all three a perfect 5/5 on cultural neutrality. Across the full dataset the judge's cultural-neutrality score is a flat $\approx 5.0$ in every condition.

\begin{table}[t]
\caption{Person-name origin by condition (\texttt{h5\_name\_origin\_crosstab.csv}).}
\label{tab:names}
\centering\small
\begin{tabular}{lcccc}
\hline
Condition & No name & Anglo/other & Hispanic & Chinese \\
\hline
Control & 366 & 21 & 1 & 0 \\
Multilingual & 307 & 16 & 72 & 0 \\
ELL-Spanish & 308 & 1 & 86 & 0 \\
ELL-Mandarin & 311 & 15 & 1 & 62 \\
\hline
\end{tabular}
\end{table}

\begin{figure}[t]
\centering
\includegraphics[width=0.8\linewidth]{figures/name_origin.png}
\caption{Person-name origin by condition. Control problems are mostly name-free; ELL labels produce strongly matched names.}
\label{fig:names}
\end{figure}

Models frequently wrote entire problems in the labeled first language unprompted (Table~\ref{tab:switch}): 53.7\% of readable ELL-Spanish problems returned in Spanish and 19.5\% of ELL-Mandarin problems in Chinese, versus near-zero switching in the control and multilingual conditions. This can be framed as helpful accommodation or as withholding the English math vocabulary an English-medium learner needs; the label alone does not tell the model which framing applies. Practically, a ``Spanish L1'' label should not be assumed to return English.

\begin{table}[t]
\caption{Output language by condition, readable items, three main families (\texttt{descriptives\_by\_condition.csv}).}
\label{tab:switch}
\centering\small
\begin{tabular}{lcccc}
\hline
Condition & $n$ & English & Spanish & Chinese \\
\hline
Control & 388 & 88.4\% & 0.0\% & 0.0\% \\
Multilingual & 395 & 89.9\% & 1.0\% & 0.0\% \\
ELL-Spanish & 395 & 38.5\% & 53.7\% & 0.0\% \\
ELL-Mandarin & 389 & 69.7\% & 0.0\% & 19.5\% \\
\hline
\end{tabular}
\end{table}

\subsection{Reliability and Judge Validity}

Reliability varied sharply by model: unreadable-JSON rates were 8.1\% (Llama), 1.5\% (Mistral), and 0.2\% (GPT-OSS); within-cell duplication was 19.3\% (Llama), 4.2\% (Mistral), and 2.4\% (GPT-OSS). At the free tier, model choice affects output quality far more than any label effect. On judge validity, correctness passed the pre-registered mark (Spearman $\rho=0.70$, 98\% agreement, $n=56$); other dimensions showed ceiling effects (both judge and human near 5.0) and are reported descriptively. In the batch's one correctness disagreement, the judge was right: a control problem stated \$20.93 where the correct total was \$21.00, which the human rater scored 5/5 and the judge caught (score 1). The human-validation batch covered only two of three families (Llama unvalidated): the least reliable model, and the one with the largest ELL trend, is also the least human-checked.

\section{Discussion}

The clearest reading is that language labels changed wording, not difficulty. Models used simpler language ($\approx 0.9$ grades) while leaving math difficulty statistically unchanged: appropriate adaptation, the opposite of the expectancy-driven harm motivating this study. Three caveats bound this reassurance. First, the null holds only for three free, mid-sized models on middle-school topics with JSON output at one temperature; it says nothing about frontier systems. Second, a null after correction is not proof of zero effect; small negative trends remain that this sample cannot rule out. Third, and most important, because models resisted lowering difficulty even for struggling students (Sect.~4.1), the design is better equipped to detect difficulty rising than falling, the exact direction the hypothesized harm would take. The null is real but asymmetric.

Model dependence is a question, not a claim: the two smaller models (Llama, Mistral) carried larger ELL trends than the larger GPT-OSS, a pattern that survived de-duplication, but with three models I cannot separate size from training recipe or provider. This motivates a pre-registered multi-model follow-up rather than a conclusion.

Unprompted language switching is the study's most open question. It can help or harm depending on facts (instructional language, student preference) the model is never told; the practical implication is that switching must be made controllable, not assumed away.

The AI judge's cultural blind spot is a load-bearing secondary finding: a judge that scores cultural neutrality at a constant $\approx 5.0$ while missing three human-flagged items is a concrete instance of the gap between aggregate judge-human agreement and specific blind spots documented elsewhere \cite{zheng2023,gallegos2024}. A high automated score should not be treated as sufficient evidence of cultural appropriateness for labeled content; human review remains necessary, not optional.

Four practical implications follow. State ability level explicitly to control difficulty, since language labels barely move it. Do not assume English output under a language-learner label. Verify mathematics with more than one checker, since a single wrong answer passed a careful human rater and was only caught by the judge. And at the free tier, prioritize model reliability (an 8\%-unreadable, 19\%-duplicate model is not interchangeable with one at 0.2\%/2\%) over any subtle fairness tuning.

Finally, a major provider's free tier returned only 138 of 540 requested Gemini responses (29 usable) over four days: independently auditing that model on a student's budget was not possible, which is itself a finding about who can audit these systems.

\section{Limitations}

The models are free, open-weight, and mid-sized (24B-120B parameters); none of the findings extend to frontier systems. The manipulation check passed only at the top of the difficulty scale (Sect.~4.1), compressing exactly the range relevant to H2's hypothesized harm; the pre-registration overstated within-model power (corrected to $\approx.67$ for $d=0.3$, Sect.~3.1); the three repeats per cell are not fully independent (addressed via clustering, Sect.~4.4); temperature was fixed at 0.7; and the JSON-output requirement may suppress behaviors visible in open-ended chat. For human validation, 56 items across two of three families were rated by a single author-rater; Llama, the least reliable model and the one with the largest ELL trend, was never human-validated, and ceiling effects restricted confirmatory validation to correctness alone. On measurement, name-origin coding (H5) is rough and after-the-fact, though the effect size is large enough to survive reasonable re-coding; the \texttt{langdetect} classifier is noisy on short text; and student labels are prompt personas, not real students, so results describe model behavior, not learning outcomes. Unreadable output (0.2-8.1\%) and duplication (2.4-19.3\%) both vary by model, and Gemini's 29 usable items support no inference. This is a single-author, single-round study; independent replication is needed.

\section{Conclusion}

Across 1,620 AI-generated problems from three free models, ELL labels did not significantly lower math difficulty after correction (pooled Cohen's $d \leq 0.25$) but did simplify language by roughly 0.9 grade levels: the models adapted the language, not the math. This finding is bounded: it applies only to mid-sized, free-tier models on middle-school topics, and the models' resistance to lowering difficulty even for struggling students means the design is better at catching difficulty rising than falling, the direction the studied harm would take. Practically, educators using AI-generated practice problems should verify mathematical correctness by hand, not assume English output under a language-learner label, and have a human, not only an AI judge, review generated content for cultural framing, since the judge in this study missed every one of the three problems a human flagged.

\section*{Data and Code Availability}

All code, data, and the pre-registration (\texttt{DESIGN.md}, filed 2026-07-03) are public at \url{https://github.com/rudrarajumc-star/prompt-sensitivity-study} (MIT license for code; CC BY 4.0 for text, figures, and tables), with a permanent version-tagged archive on Zenodo (DOI: \href{https://doi.org/10.5281/zenodo.21824157}{10.5281/zenodo.21824157}). The pipeline (\texttt{build\_prompt\_matrix.py} $\to$ \texttt{generate.py} $\to$ \texttt{evaluate.py} $\to$ \texttt{judge.py} $\to$ \texttt{analyze.py}) regenerates every table and figure from raw data; figures regenerate at 300 dpi via \texttt{make\_report\_figures.py}. Every quantitative claim in this paper traces to a file in the repository. A longer version of this paper, with full appendices (complete prompt wordings, exact model identifiers, additional selected generations, and an extended reproducibility log), is available in the GitHub repository above and archived at the Zenodo DOI above; this submission is the condensed, self-contained version of that work.

\section*{Ethics Statement}

This study used no human participants or student data; all ``students'' are fictional prompt personas. I am the only human contributor, and I provided the blinded ratings of model output myself. The audited fairness risk, automatically lowering mathematical expectations when a student's language background is mentioned, targets a protected, often under-served population; the audit's purpose is protective, and I report results in full, including the reassuring null, its bounds, and the unresolved questions. All protocol amendments are disclosed with timestamps (Sect.~3.4); data and code are fully released for independent verification.

\begin{credits}
\subsubsection{\ackname}
I thank Prof. Parameswari Krishnamurthy (IIIT Hyderabad, Language Technologies Research Centre) for a research discussion that informed the future-work directions in Sect.~5; she bears no responsibility for the analysis or any errors. I received no funding; all model access used free, no-credit-card provider tiers.

AI tool use is disclosed in full in Sect.~3.3.

\subsubsection{\discintname}
I declare no competing interests.
\end{credits}

\begin{thebibliography}{22}

\bibitem{blodgett2020} Blodgett, S.L., Barocas, S., Daum\'e III, H., Wallach, H.: Language (technology) is power: A critical survey of ``bias'' in NLP. In: Proceedings of ACL 2020. pp. 5454--5476 (2020)

\bibitem{christ2024} Christ, B., Kropko, J., Hartvigsen, T.: MATHWELL: Generating educational math word problems using teacher annotations. In: Findings of EMNLP 2024 (2024)

\bibitem{cobbe2021} Cobbe, K., Kosaraju, V., Bavarian, M., et al.: Training verifiers to solve math word problems. arXiv:2110.14168 (2021)

\bibitem{cohen1988} Cohen, J.: Statistical Power Analysis for the Behavioral Sciences. 2nd edn. Lawrence Erlbaum Associates (1988)

\bibitem{deshpande2023} Deshpande, A., Murahari, V., Rajpurohit, T., Kalyan, A., Narasimhan, K.: Toxicity in ChatGPT: Analyzing persona-assigned language models. In: Findings of EMNLP 2023 (2023)

\bibitem{gallegos2024} Gallegos, I.O., Rossi, R.A., Barrow, J., et al.: Bias and fairness in large language models: A survey. Computational Linguistics 50(3), 1097--1179 (2024)

\bibitem{gupta2023} Gupta, S., Shrivastava, V., Deshpande, A., et al.: Bias runs deep: Implicit reasoning biases in persona-assigned LLMs. In: ICLR 2024 (2023)

\bibitem{holm1979} Holm, S.: A simple sequentially rejective multiple test procedure. Scandinavian Journal of Statistics 6(2), 65--70 (1979)

\bibitem{kasneci2023} Kasneci, E., Sessler, K., K\"uchemann, S., et al.: ChatGPT for good? On opportunities and challenges of large language models for education. Learning and Individual Differences 103, 102274 (2023)

\bibitem{kincaid1975} Kincaid, J.P., Fishburne, R.P., Rogers, R.L., Chissom, B.S.: Derivation of New Readability Formulas for Navy Enlisted Personnel. Research Branch Report 8-75, Naval Technical Training Command (1975)

\bibitem{liang2022} Liang, P., Bommasani, R., Lee, T., et al.: Holistic evaluation of language models. arXiv:2211.09110 (2022)

\bibitem{mackinnon1985} MacKinnon, J.G., White, H.: Some heteroskedasticity-consistent covariance matrix estimators with improved finite sample properties. Journal of Econometrics 29(3), 305--325 (1985)

\bibitem{nakatani2010} Nakatani, S.: Language detection library for Java (2010), Python port: langdetect

\bibitem{oakes1985} Oakes, J.: Keeping Track: How Schools Structure Inequality. Yale University Press (1985)

\bibitem{rosenthal1968} Rosenthal, R., Jacobson, L.: Pygmalion in the Classroom: Teacher Expectation and Pupils' Intellectual Development. Holt, Rinehart \& Winston (1968)

\bibitem{sclar2023} Sclar, M., Choi, Y., Tsvetkov, Y., Suhr, A.: Quantifying language models' sensitivity to spurious features in prompt design. In: ICLR 2024 (2023)

\bibitem{sheng2019} Sheng, E., Chang, K.W., Natarajan, P., Peng, N.: The woman worked as a babysitter: On biases in language generation. In: Proceedings of EMNLP-IJCNLP 2019. pp. 3407--3412 (2019)

\bibitem{shi2022} Shi, F., Suzgun, M., Freitag, M., et al.: Language models are multilingual chain-of-thought reasoners. In: ICLR 2023 (2022)

\bibitem{speer2022} Speer, R.: rspeer/wordfreq: v3.0. Zenodo (2022). \url{https://doi.org/10.5281/zenodo.7199437}

\bibitem{steele1995} Steele, C.M., Aronson, J.: Stereotype threat and the intellectual test performance of African Americans. Journal of Personality and Social Psychology 69(5), 797--811 (1995)

\bibitem{umansky2021} Umansky, I.M., Dumont, H.: English learner labeling: How English learner classification in kindergarten shapes teacher perceptions of student skills. American Educational Research Journal 58(5), 993--1031 (2021)

\bibitem{zheng2023} Zheng, L., Chiang, W.L., Sheng, Y., et al.: Judging LLM-as-a-judge with MT-Bench and Chatbot Arena. In: NeurIPS 2023 Datasets and Benchmarks (2023)

\end{thebibliography}

\end{document}
